\documentclass[11pt, twocolumn]{article}
\usepackage{graphicx}
\usepackage{multicol}
\usepackage{tabularx}
\usepackage[margin=1in]{geometry}
\usepackage{url}
\usepackage{listings}
%\usepackage[style=ieee]{biblatex}

\renewcommand{\abstractname}{Resumen}
\renewcommand{\refname}{Referencias}

\begin{document}

\begin{titlepage}
	\centering
	\LARGE
	\textbf{INSTITUTO TECNOLÓGICO Y DE ESTUDIOS SUPERIORES DE OCCIDENTE}
	
	\large
	\textbf{Departamento de Electrónica, Sistemas e Informática}\\
	\textbf{Maestría en Sistemas Computacionales}\\
	\vspace{1.0cm}
	\textbf{Procesamiento digital de señales}
	\textbf{Otoño 2022}
	\includegraphics[scale=1.2]{images/iteso_logo}

	\LARGE
	\textbf{Tarea 0: Resumen de artículos sobre Procesamiento de Señales}
	\vfill
	
	\textsf{741645, erich.cordoba@iteso.mx}\\
	\textsc{Tlaquepaque, Jalisco. Agosto 2022}
\end{titlepage}


\begin{abstract}
Realizar un resumen seleccionando dos artículos propuestos y un tercero a criterio personal.

\end{abstract}

\section{Introducción}
Hablar sobre los papers que se eligieron. Probando \cite{haykin}, tambien \cite{golden_age} y finalmente \cite{dps_dead}.

\section{Primer artículo}
Descripción del primer artíclo

\section{Segundo artículo}
Hablar sobre el segundo artículo

\section{Tercer articulo}
Hablar sobre el tercer artículo


\section{Referencias}

\bibliographystyle{ieeetr}
\bibliography{referencias}

\end{document}